\begin{singlespace}
\setlength{\parindent}{0pt}
\raggedright % ABNT NBR 6023: referências alinhadas à esquerda
\setlength{\parskip}{0.6em}

BOURAHMOUNE, K.; AMAGASA, T. AI-powered posture training: application of machine
learning in sitting posture recognition using the LifeChair smart cushion. In:
INTERNATIONAL JOINT CONFERENCE ON ARTIFICIAL INTELLIGENCE, 28., 2019, Macau.
\emph{Proceedings} [...]. [S.l.]: IJCAI, 2019. p. 5808--5814. DOI:
10.24963/ijcai.2019/805.

BREEN, P. P.; NISAR, A.; ÓLAIGHIN, G. Evaluation of a single accelerometer based
biofeedback system for real-time correction of neck posture in computer users.
In: ANNUAL INTERNATIONAL CONFERENCE OF THE IEEE ENGINEERING IN MEDICINE AND
BIOLOGY SOCIETY, 31., 2009, Minneapolis. \emph{Proceedings} [...]. Piscataway:
IEEE, 2009. p. 7269--7272. DOI: 10.1109/IEMBS.2009.5334726.

DUL, J.; WEERDMEESTER, B. \emph{Ergonomia prática}. 3. ed. São Paulo: Edgard
Blücher, 2012.

FERREIRA, J. A.; SOUZA, R. A. de; CUBAS, J. J. M. Desenvolvimento de um sistema
de monitoramento e alerta de postura utilizando sensores integrados em cadeira:
aplicação do BackWatch para melhoria da ergonomia. In: FEIRA BRASILEIRA DE
CIÊNCIAS E ENGENHARIA (FEBRACE), 2025, São Paulo. \emph{Mostra virtual}:
projeto SAU-8347. São Paulo: FEBRACE, 2025. Pôster. Disponível em: \url{https://virtual.febrace.org.br/2025/SAU/8347}. Acesso em: 8 nov. 2026.

GRANDJEAN, E. \emph{Manual de ergonomia}: adaptando o trabalho ao homem. 4. ed.
Porto Alegre: Bookman, 1998.

HALLER, M. et al. Finding the right way for interrupting people improving their
sitting posture. In: IFIP TC 13 INTERNATIONAL CONFERENCE ON HUMAN-COMPUTER
INTERACTION (INTERACT), 13., 2011, Lisboa. \emph{Human-Computer Interaction ---
INTERACT 2011}: part II. Berlin: Springer, 2011. p. 1--17. (Lecture Notes in
Computer Science, v. 6947). DOI: 10.1007/978-3-642-23771-3\_1.

IIDA, I.; GUIMARÃES, L. B. M. \emph{Ergonomia}: projeto e produção. 3. ed. São
Paulo: Edgard Blücher, 2016.

INTERLINK ELECTRONICS. \emph{FSR\textsuperscript{\textregistered} 400 Series data
sheet}. P/N PDS-10004-C. [S.l.]: Interlink Electronics, [2012?a]. Disponível em: \url{https://www.teachengineering.org/content/uconn_/activities/uconn-2662-contractions-calculations-force-resistance-activity/uconn-2662-contractions-calculations-FSR400series-datasheet.pdf}.
Acesso em: 8 nov. 2026.

INTERLINK ELECTRONICS. \emph{FSR\textsuperscript{\textregistered} integration
guide}. Doc. 94-00004 Rev. C. [S.l.]: Interlink Electronics, [2012?b].
Disponível em: \url{https://www.pololu.com/file/0J749/FSR400-Series-Integration-Guide-13.pdf}.
Acesso em: 8 nov. 2026.

INTERNATIONAL ORGANIZATION FOR STANDARDIZATION. \emph{ISO 11226}: ergonomics ---
evaluation of static working postures. Geneva: ISO, 2000.

KAPPATTANAVAR, A. M. et al. Monitoring of sitting postures with sensor networks in
controlled and free-living environments: systematic review. \emph{JMIR Biomedical
Engineering}, v. 6, n. 1, e21105, 2021. DOI: 10.2196/21105.

KUO, Y.-L. et al. Sitting posture during prolonged computer typing with and
without a wearable biofeedback sensor. \emph{International Journal of
Environmental Research and Public Health}, v. 18, n. 10, 5430, 2021. DOI:
10.3390/ijerph18105430.

KUORINKA, I. et al. Standardised Nordic questionnaires for the analysis of
musculoskeletal symptoms. \emph{Applied Ergonomics}, v. 18, n. 3, p. 233--237,
1987.

LUINGE, H. J.; VELTINK, P. H. Measuring orientation of human body segments using
miniature gyroscopes and accelerometers. \emph{Medical and Biological Engineering
and Computing}, v. 43, p. 273--282, 2005. DOI: 10.1007/BF02345966.

MCATAMNEY, L.; CORLETT, E. N. RULA: a survey method for the investigation of
work-related upper limb disorders. \emph{Applied Ergonomics}, v. 24, n. 2,
p. 91--99, 1993.

MEYER, J. et al. Design and modeling of a textile pressure sensor for sitting
posture classification. \emph{IEEE Sensors Journal}, v. 10, n. 8,
p. 1391--1398, 2010. DOI: 10.1109/JSEN.2009.2037330.

MOHAMMADIAN et al. Musculoskeletal disorders among office workers: prevalence,
ergonomic risk factors, and their interrelationships. \emph{Scientific Reports},
v. 15, 2025. DOI: 10.1038/s41598-025-30155-6.

MUTLU, B. et al. Robust, low-cost, non-intrusive sensing and recognition of
seated postures. In: ANNUAL ACM SYMPOSIUM ON USER INTERFACE SOFTWARE AND
TECHNOLOGY, 20., 2007, Newport. \emph{Proceedings} [...]. New York: ACM, 2007.
p. 149--158. DOI: 10.1145/1294211.1294237.

OKEZUE, O. C. et al. Work-Related Musculoskeletal Disorders among Office Workers
in Higher Education Institutions: A Cross-Sectional Study. \emph{Ethiopian
Journal of Health Sciences}, v. 30, n. 5, p. 715--724, 2020.

PAREDES-MADRID, L. et al. Underlying physics of conductive polymer composites and
force sensing resistors (FSRs): a study on creep response and dynamic loading.
\emph{Materials}, v. 10, n. 11, 1334, 2017. DOI: 10.3390/ma10111334.

PEDLEY, M. \emph{Tilt sensing using a three-axis accelerometer}. Application
Note AN3461, Rev. 6. [S.l.]: Freescale Semiconductor, 2013. Disponível em: \url{https://www.nxp.jp/docs/en/application-note/AN3461.pdf}. Acesso em: 8 nov. 2026.

RIBEIRO, D. C.; SOLE, G.; ABBOTT, J. H.; MILOSAVLJEVIC, S. The effectiveness of a
lumbopelvic monitor and feedback device to change postural behavior: a
feasibility randomized controlled trial. \emph{Journal of Orthopaedic \& Sports
Physical Therapy}, v. 44, n. 9, p. 702--711, 2014. DOI: 10.2519/jospt.2014.5009.

RIBEIRO, D. C. et al. Effectiveness of a lumbopelvic monitor and feedback device
to change postural behaviour: the ELF cluster randomised controlled trial.
\emph{Occupational and Environmental Medicine}, v. 77, n. 7, p. 462--469, 2020.
DOI: 10.1136/oemed-2019-106293.

ROBERTSON, M.; CIRIELLO, V.; GARABET, A. Office ergonomics training and a
sit-stand workstation: effects on musculoskeletal and visual symptoms and
performance of office workers. \emph{Applied Ergonomics}, v. 44, n. 1,
p. 73--85, 2013.

ROH, J. et al. Sitting posture monitoring system based on a low-cost load cell
using machine learning. \emph{Sensors}, v. 18, n. 1, 208, 2018. DOI:
10.3390/s18010208.

SONNE, M.; VILLALTA, D. L.; ANDREWS, D. M. Development and evaluation of an
office ergonomic risk checklist: ROSA --- Rapid Office Strain Assessment.
\emph{Applied Ergonomics}, v. 43, n. 1, p. 98--108, 2012.

TAN, H. Z.; SLIVOVSKY, L. A.; PENTLAND, A. A sensing chair using pressure
distribution sensors. \emph{IEEE/ASME Transactions on Mechatronics}, v. 6, n. 3,
p. 261--268, 2001. DOI: 10.1109/3516.951364.

WONG, W. Y.; WONG, M. S. Trunk posture monitoring with inertial sensors.
\emph{European Spine Journal}, v. 17, n. 5, p. 743--753, 2008. DOI:
10.1007/s00586-008-0586-0.

XU, W. et al. eCushion: a textile pressure sensor array design and calibration
for sitting posture analysis. \emph{IEEE Sensors Journal}, v. 13, n. 10,
p. 3926--3934, 2013. DOI: 10.1109/JSEN.2013.2259589.

ZEMP, R. et al. Application of machine learning approaches for classifying
sitting posture based on force and acceleration sensors. \emph{BioMed Research
International}, v. 2016, 5978489, 2016. DOI: 10.1155/2016/5978489.

ZHENG, Y.; MORRELL, J. B. Comparison of visual and vibrotactile feedback methods
for seated posture guidance. \emph{IEEE Transactions on Haptics}, v. 6, n. 1,
p. 13--23, 2013. DOI: 10.1109/TOH.2012.3.
\end{singlespace}
